\section{From Forecast Availability to Realized Gain}
We test whether access to learned paths adds value over a selector
that can already choose among five causal motion alternatives.
The all-recipient hindsight comparison fixes the CV fallback,
per-query intervention caps, the 2\% risk rule, and the OLS4
accuracy reference. The motion bank has positive lower utility on
\CausalPositive{} of 102 recording-seed views. Adding the three
learned forecasts raises this to \MixedPositive{} views, with the
new positives all from one recording. The prior shared-motion
transform adds no further positive views. The source-balanced
increase is \CapacityGain; the solver uses labels and this is
not evidence of a past-only policy improvement.

In the past-only comparison, the two joint policies make
\MenuSwitches{} genuine changes each, with identical counts per
query. The mixed bank improves available-label utility by
\MenuObserved, while its conservative lower contrast is \MenuLower.
Observed relative-pair ADE reduction is \MenuPair{} pixels.
The two predicted screens share \CommonScreenSwitches{} accepted
switches, all identical motion choices; the intersection is not
an alternative operating policy and does not certify risk.

\paragraph{Learning the neural increment directly.}
Rather than estimate benefit from the difference of two separate
CV-relative scores, we learn the change from the actual selected
motion forecast. We fix the recipients, clip both incremental
scores to the same pairwise disagreement range, and change a path
only when its predicted increment is positive. For each held
source, anchors on the other 13 sources are constructed using
fits that also exclude their own source. The ridge heads see
the same past descriptor family centered on the selected
motion path. Their target is its available-label gain minus
maximum paired disagreement on wholly unobserved rows, divided
by the past motion scale. These changes do not alter the CV risk
reference or any evaluation rule.

At the same recipients and counts, the direct score changes the
conservative lower utility by \DirectVsCV{} relative to the
CV-difference control and by \DirectVsAnchor{} relative to
retaining the motion anchors. The respective whole-recording
pass counts (out of 102) are \DirectPass{} for the direct
score, \DifferencePass{} for CV-difference, and \AnchorPass{}
for the motion anchor. The direct-minus-CV-difference contrast
on observed labels is \DirectKnown; the lower envelope subtracts
another \DirectMissing{} points for wholly unobserved rows.
The unchanged predicted screen retains only \DirectScreenSwitches{}
direct-score switches and passes \DirectScreenPass{} of 102 views.
The observed relative-pair ADE reduction
for direct versus anchor is \DirectPair{} pixels. These results
are development-only. Learning a different score does not
by itself convert a label-aware witness into a risk-supported
decision on a new recording.
